\rok{Март 2026.}{Б}

Трећи поправни тест из Алгоритама и структура података

Први практични тест одржан 06.03.2026.

\zadatak{1}{Имплементација DFSVisit методе \textit{Depth-First-Search} алгоритма}
Написати имплементацију \texttt{DFSVisit} методе истоименог алгоритма.

\textbf{Додатне напомене за решавање задатка:}
\begin{itemize}
	\item Написати алгоритам.
	\item У \texttt{Main} методи креирати произвољан граф. Обезбедити начин за тестирање и проверу нове функционалности, односно позив \texttt{DFS} методе.
	\item У \textit{template}-у је метода \texttt{private void DFSVisit(LinkedList.Node node, int nodeIndex)} која је основна метода за рекурзивно извршавање DFS алгоритма.
\end{itemize}

\textbf{Потпис методе:}
\begin{Verbatim}[fontsize=\small]
private void DFSVisit(LinkedList.Node node, int nodeIndex)
\end{Verbatim}

\zadatak{2}{Чишћење дупликата (Keep First Occurrence)}
Написати методу која из прослеђеног стринга уклања све накнадне појаве карактера који су се већ појавили, остављајући само њихову прву појаву. Резултат исписати у конзоли. Алгоритам треба да буде временске комплексности $O(n)$ и \textit{case-insensitive}.

\textbf{Додатни захтеви за решавање задатка:}
\begin{itemize}
	\item Користити \textit{Hash} табелу из шаблона за чување броја појављивања сваког карактера.
	\item Имплементацију писати у методи: \texttt{public static void keepFirstOccurrence (string str)}.
	\item Временска комплексност треба да буде $O(n)$.
	\item Игнорисати разлику између малих и великих слова.
\end{itemize}

\textbf{Потпис методе:}
\begin{Verbatim}[fontsize=\small]
public static void keepFirstOccurrence (string str)
\end{Verbatim}

\textbf{Пример:}
\begin{itemize}
	\item \textbf{Улаз 1:} \texttt{str = "Programiranje"}
	\item \textbf{Излаз 1:} \texttt{"Progaminje"} (карактери \texttt{'r'} и \texttt{'a'} се понављају)
\end{itemize}

\zadatak{3}{Бројање чворова на специфичном нивоу N}
У оквиру класе \texttt{BinarySearchTree} написати методу која броји колико се укупно чворова налази на тачно одређеном нивоу (дубини) \texttt{N}, где је корен стабла на нивоу 0.

\textbf{Додатни захтеви за решавање задатка:}
\begin{itemize}
	\item Јавна метода: \texttt{public int countNodesAtLevel (int targetLevel)}.
	\item Приватна метода: \texttt{private int countNodesAtLevel (Node current, int currentLevel, int targetLevel)}.
	\item Подразумева се да стабло има барем два чвора.
	\item Алгоритам \textbf{ОБАВЕЗНО} писати у оквиру класе \texttt{BinarySearchTree}.
\end{itemize}

\textbf{Потписи метода:}
\begin{Verbatim}[fontsize=\small]
public int countNodesAtLevel (int targetLevel)
private int countNodesAtLevel (Node current, int currentLevel, int targetLevel)
\end{Verbatim}

\textbf{Примери:}

\begin{center}
\begin{minipage}{\textwidth}
\centering
\begin{forest}
for tree={
	circle,
	draw,
	thick,
	fill=gray!20,
	minimum size=8mm,
	inner sep=1pt,
	s sep=8mm,
	l sep=12mm,
	edge={-, thick},
	font=\small
}
[8
	[4
		[2
			[, phantom]
			[3]
		]
		[5]
	]
	[9
		[, phantom]
		[11]
	]
]
\end{forest}

\vspace{0.3cm}
\textbf{Слика 1.} Пример BST-а.
\end{minipage}
\end{center}

\begin{itemize}
	\item \textbf{Улаз 1:} \texttt{targetLevel = 3}
	\item \textbf{Излаз 1:} 1 (\texttt{"3"} се налази на тој дубини)
\end{itemize}

\begin{center}
\begin{minipage}{\textwidth}
\centering
\begin{forest}
for tree={
	circle,
	draw,
	thick,
	fill=gray!20,
	minimum size=8mm,
	inner sep=1pt,
	s sep=8mm,
	l sep=12mm,
	edge={-, thick},
	font=\small
}
[10
	[4
		[2
			[, phantom]
			[3]
		]
		[6
			[5]
			[9]
		]
	]
	[20
		[12
			[, phantom]
			[14]
		]
		[29
			[27]
			[, phantom]
		]
	]
]
\end{forest}

\vspace{0.3cm}
\textbf{Слика 2.} Пример BST-а.
\end{minipage}
\end{center}

\begin{itemize}
	\item \textbf{Улаз 2:} \texttt{targetLevel = 2}
	\item \textbf{Излаз 2:} 4 (\texttt{"2"}, \texttt{"6"}, \texttt{"12"} и \texttt{"29"})
\end{itemize}

\sablon{Шаблон другог поправног теста март 2026. група Б}{2025/Mart/GrupaB/AISP2025_PopravniTest3_Test2_GrupaB.linq}
